\documentclass{l4dc2026}
\usepackage{iftex,booktabs}
\ifPDFTeX
\usepackage{times}
\else
\usepackage{fontspec}
\setmainfont{Times New Roman}
\fi
\jmlrvolume{}
\jmlrproceedings{ESE 6180}{ESE 6180 Results Note}
\jmlrworkshop{ESE 6180 Results Note}
\title[History, Coverage, and Chunking]{History and Chunking from Two to Sixteen States}
\author{\Name{Nihar Ballamudi}\and\Name{Kyle Horton}\\\addr University of Pennsylvania}
\begin{document}
\maketitle
\begin{abstract}
We compare the original Simchowitz two-state example with coupled systems of 2, 4, 8, and 16 states. Noisy training actions provide corrective examples; with enough history, single-action control nearly matches a known-model observer. All testing is noiseless.
\end{abstract}

\section{The experiment}
\citet{simchowitz2025} give two systems with experts $K_1,K_2$ that agree on every demonstrated state but respond differently outside those states. Each expert stabilizes its own system; the wrong expert can destabilize the other. This shows why matching demonstrated actions does not guarantee a good controller.

We retain the paper's original $n=2$ pair at $\mu=1/8$ and add a separate coupled family for $n=2,4,8,16$. The new family is not made of repeated two-state blocks. Each system has one feedback loop through all its coordinates:
\[
x_1^+=1.02x_1+0.04x_n+u_1,\quad
x_2^+=-x_1+u_2,\quad x_j^+=x_{j-1}+u_j\ (j\ge3).
\]
Only $x_n$ is observed. Its second system replaces the first column of the state matrix by $-0.995e_1$. Set $K_i=F-A_i$, with $F=\operatorname{diag}(0,\ldots,0,0.75)$. Both experts produce identical demonstrations when only $x_n$ is nonzero, but cross-applying them gives an unstable eigenvalue $+2.015$ or $-2.015$. Exact reconstruction requires $n$ observations. The original-paper $n=2$ and coupled-family $n=2$ are different systems and are labeled separately.

We compare single actions and 16-action predictions, each with and without history. History $r$ contains the current and past observations plus the intervening applied actions. Without history, $r=1$. Single-action control uses prediction horizon $H=1$ and execution horizon $\ell=1$. A 16-action predictor executes $\ell=1,2,4,8,16$ actions before observing and predicting again. Histories cover 1 through $n+1$, plus shared lengths 2, 4, 8, and 16.

\section{The training and test data}
Noiseless training data follow the expert and never leave the demonstrated line. Noisy training data add independent Gaussian changes of standard deviation 0.05 to the applied actions. At each resulting state, the training answer is the exact action sequence the expert would take without further changes. Thus noise changes the states visited; it does not corrupt the expert answers. Inputs retain the actions actually applied.

Ridge regression fits a linear action rule with coefficient penalty $10^{-6}$. Five seeds each generate 128 trajectories of 100 steps per system and dataset. Target times and 16-step expert initialization match across dimensions. Executions share the same predictor; applying only its first action matches a separately trained single-action rule.

Testing uses 40 shared visible starting amplitudes in $[-1,1]$. After filling history, we introduce a hidden-state change with total magnitude 0.0005--0.0015, either in $x_1$ or spread across hidden coordinates. These are the same hidden axis at $n=2$, not independent spatial directions. The learner then acts independently for 200 steps, with no further noise or expert queries. Error averages squared full-state distance from a separate expert rollout over time, both change types, and training seeds. We separately check eventual growth or decay and selected 1000-step examples.

\section{The results}
\begin{center}\small
\begin{tabular}{ll lrrrrr}\toprule
System & $n$ & Training data & \shortstack{History\\$r$} & \shortstack{Prediction\\horizon $H$} & \shortstack{Execution\\horizon $\ell$} & Mean test MSE & Stable seeds\\\midrule
Original & 2 & Noiseless & 2 & 1 & 1 & $6.09\times10^{9}$ & 0/5\\
Original & 2 & Noisy & 2 & 1 & 1 & $7.07\times10^{-9}$ & 5/5\\
Coupled & 2 & Noiseless & 2 & 1 & 1 & $5.04\times10^{-6}$ & 5/5\\
Coupled & 2 & Noisy & 2 & 1 & 1 & $1.75\times10^{-8}$ & 5/5\\
Coupled & 2 & Noisy & 2 & 16 & 16 & $1.36\times10^{-7}$ & 5/5\\
Coupled & 4 & Noisy & 4 & 1 & 1 & $4.51\times10^{-8}$ & 5/5\\
Coupled & 4 & Noisy & 4 & 16 & 16 & $2.05\times10^{-7}$ & 5/5\\
Coupled & 8 & Noisy & 8 & 1 & 1 & $1.49\times10^{-7}$ & 5/5\\
Coupled & 8 & Noisy & 8 & 16 & 16 & $3.90\times10^{-7}$ & 5/5\\
Coupled & 16 & Noisy & 16 & 1 & 1 & $5.47\times10^{-7}$ & 5/5\\
Coupled & 16 & Noisy & 16 & 16 & 16 & $5.54\times10^{-7}$ & 5/5\\
\bottomrule\end{tabular}
\end{center}

All systems share seeds, sample counts, and timing; both $n=2$ systems also share test states. Dimension-scaling plots use coupled $n=2$; a separate graph compares it with the original pair.

With noisy training data, single-action control at history $r=n$ is stable for every seed and nearly matches the known-model observer. Coupled-family errors are $1.75\times10^{-8}$, $4.51\times10^{-8}$, $1.49\times10^{-7}$, and $5.47\times10^{-7}$. At $r=n-1$, all five single-action fits are unstable at every dimension. More state variables therefore increase the history needed to learn accurate corrections in these runs.

Executing sixteen actions at $r=n$ increases mean error by factors 7.82, 4.56, 2.63, and 1.01 across the four dimensions. It delays useful corrections rather than improving the learned rule. With noiseless data, history can instead stabilize a controller without teaching the missing corrections: the coupled $n=2$ fit at history two is stable but has about 289 times the noisy-data error. The original pair's noiseless fit remains unstable at history two. These differences are why the two $n=2$ systems must not be pooled.

The findings concern known action perturbations during training, one initial test departure, and no continuing test noise. Stability and tracking accuracy are separate: a small finite-rollout error is not proof of eventual decay.

\footnotesize
\begin{thebibliography}{1}
\bibitem[Simchowitz et al.(2025)]{simchowitz2025}
M. Simchowitz, D. Pfrommer, and A. Jadbabaie.
The pitfalls of imitation learning when actions are continuous.
\emph{arXiv:2503.09722}, 2025.
\end{thebibliography}
\end{document}
